% Figure 25.2 : un volume CSI en WaitForFirstConsumer, du Pod au montage (événements réels du chapitre 25).
\documentclass[tikz,border=4pt]{standalone}
\usepackage{quai}
\begin{document}
\begin{tikzpicture}[x=1.12cm,y=1cm,
    acteur/.style={qbox, font=\ttfamily\scriptsize, minimum width=2.4cm, minimum height=0.8cm, inner sep=2pt, align=center},
    m/.style={qarr},
    mt/.style={font=\scriptsize, fill=QPaper, inner sep=1.2pt}]
  \node[qcadre, dashed, draw=QMuted, minimum width=6.4cm, minimum height=1.35cm] at (10.5,0.12) {};
  \node[qlbl, font=\scriptsize, anchor=south] at (10.5,0.8) {Pod csi-hostpathplugin};
  \foreach \x/\n/\t/\c in {0/kc/kubectl/QGris, 3/api/API server\\+ etcd/QBleu, 6/sc/scheduler/QSarcelle, 9/pr/csi-provisioner/QSarcelle, 12/pl/pilote\\hostpath/QOcre, 15/kl/kubelet/QVert} {
    \node[acteur, draw=\c, fill=\c Bg] (\n) at (\x,0) {\t};
    \draw[QMuted, line width=0.5pt, dashed] (\x,-0.45) -- (\x,-9.5);
  }
  \newcommand{\msg}[5]{%
    \draw[m] (#1,#3) -- node[mt, above] {#4} (#2,#3);
    \node[qnum=QBleu] at ({#1+0.45*sign(#2-#1)},#3-0.32) {\scriptsize #5};}
  \msg{0}{3}{-1.0}{PVC journal, puis Pod journal}{1}
  \msg{3}{6}{-1.9}{Pod à placer}{2}
  \msg{6}{3}{-2.8}{selected-node=minikube}{3}
  \msg{3}{9}{-3.7}{PVC à provisionner}{4}
  \msg{9}{12}{-4.6}{CreateVolume}{5}
  \msg{9}{3}{-5.5}{crée le PV, lié à la PVC}{6}
  \msg{6}{3}{-6.4}{nodeName = minikube}{7}
  \msg{3}{15}{-7.3}{Pod pour ce nœud}{8}
  \msg{15}{12}{-8.2}{NodePublishVolume}{9}
  \node[qlbl, font=\scriptsize, anchor=west, align=left] at (15.25,-8.9) {montage dans\\{\ttfamily /var/lib/kubelet/pods/}\\{\ttfamily <uid>/volumes/}\\{\ttfamily kubernetes.io\textasciitilde csi/}, puis\\démarrage du conteneur};
  \node[qlbl, font=\scriptsize, anchor=north, align=center] at (10.5,-9.6) {gRPC par la socket {\ttfamily /var/lib/kubelet/plugins/csi-hostpath/csi.sock}};
\end{tikzpicture}
\end{document}
